\section{PatchTST: Patches as Tokens for Time Series}

\subsection{Motivation}

Standard Transformer self-attention has complexity $O(T^2 d)$ where $T$
is the sequence length. For price forecasting with monthly lags, $T=720$
(30 days) gives 518,400 attention operations per head. For 56-day context
($T=1344$), this is $1{,}806{,}336$. On a single GPU or MPS, this is the
main reason TFT cannot go beyond ctx=2016 economically.

PatchTST (Nie et al., NeurIPS 2023) fixes this by treating contiguous
subsequences as tokens. A patch of length $P$ reduces the effective
sequence length to $\lceil T/P \rceil$, so attention cost drops from
$O(T^2)$ to $O((T/P)^2)$. For $P=24$ (one day per patch) and $T=1344$,
the effective sequence is $56$ tokens instead of $1344$.

\subsection{Architecture}

\begin{enumerate}
  \item \textbf{Patching}: the input series is split into overlapping
        patches of length $P$ with stride $S$.
        Total patches: $n\_\text{patches} = (T - P) / S + 1$.
        Example: $T=1344$, $P=24$, $S=12$ $\to$ $n = 110$ patches.
  \item \textbf{Projection}: each patch is linearly projected to a
        $d\_\text{model}$-dimensional embedding.
  \item \textbf{Positional encoding}: sinusoidal encoding added to
        patch embeddings (patch index, not hour index).
  \item \textbf{Multi-head self-attention}: over the $n$ patch tokens.
        Cost: $O(n^2 d)$ instead of $O(T^2 d)$.
  \item \textbf{Channel independence}: each variate (price, solar,
        wind, TSO) is encoded \emph{independently}. No cross-variate
        attention. This is both a limitation and a strength --- it
        reduces parameter count but loses variate interactions.
  \item \textbf{Output}: flatten the patch representations, project
        to 24-hour quantile forecasts.
\end{enumerate}

\subsection{Patch length vs context tradeoffs}

\begin{center}
\begin{tabular}{lccc}
\toprule
patch\_len & intraday signal & context at same cost & merit-order capture \\
\midrule
12h & strong & medium & morning/evening splits \\
24h & medium & high & whole-day pattern \\
48h & weak & very high & weekday vs weekend \\
\bottomrule
\end{tabular}
\end{center}

For the price task, the merit order operates at the hourly level:
solar suppresses hours 10--15, gas peaks hours 18--21.
A 24h patch collapses both into one token.
A 12h patch keeps the AM/PM distinction.

\subsection{Channel independence: the tradeoff}

PatchTST (original version) treats each feature independently.
When predicting price from solar, wind, and TSO:
\begin{itemize}
  \item Solar's encoder never ``sees'' wind's encoder.
  \item The price predictor gets concatenated embeddings, not
        attended cross-variate representations.
\end{itemize}

This hurts when the signal is an \emph{interaction} (e.g.\ low wind
AND high demand $\to$ spike). TFT's temporal self-attention operates
across all features jointly; its attention heads can learn this
cross-feature dependency.

The PatchTST-MLP variant (also in the paper) adds a cross-variate
mixing step. We did not implement this.

\subsection{Why PatchTST might beat TFT on this task}

\begin{itemize}
  \item \textbf{Fewer parameters}: PatchTST at d=64 has $\sim$50k params
        vs TFT's 1.27M. At 300--400 training samples, smaller wins.
  \item \textbf{Longer context at the same cost}: ctx=2016h with
        P=24 gives 84 tokens vs 2016 tokens for TFT. Attention is
        much cheaper; training is faster; longer history fits.
  \item \textbf{Simpler training}: no VSN, no LSTM encoder, no
        quantile-gating. The architecture is a more direct map from
        patches to output.
\end{itemize}

\subsection{Why PatchTST might lose}

\begin{itemize}
  \item Channel independence loses cross-variate interactions.
        The solar $\times$ demand interaction (duck curve) requires
        joint encoding.
  \item Patch compression loses intraday signal at $P=24$ or $P=48$.
        The merit order is an \emph{hourly} effect; a daily patch
        averages it out.
  \item The price series is non-stationary in a regime-switching way
        (gas crisis, solar growth). Patch tokens are not explicitly
        normalised per-token (unlike TFT's instance normalisation).
\end{itemize}

\subsection{Expected result from the sweep}

Based on theoretical analysis:
\begin{itemize}
  \item Best patch length: 12h (preserves intraday signal, still cheaper
        than raw attention at ctx=1344).
  \item Best stride: 6h (50\% overlap, more patches, richer context).
  \item Best context: 1344h or 2016h (regime memory is the gap vs tabular).
  \item Expected val pinball: 0.120--0.130 range (vs TFT 0.1157).
        If PatchTST matches TFT on screening, walk-forward wins likely
        due to fewer parameters.
\end{itemize}

\textit{Fill in actual results after run\_patchtst\_sweep.py completes.}

\subsection{Production bug: zero-variance columns blow up val loss}

This is a real ML-ops story worth knowing.

\textbf{Symptom.} First PatchTST sweep showed val pinball of $879$ vs training
pinball of $0.26$ --- a 3000$\times$ gap, starting at epoch 0 even before
any useful learning had occurred.

\textbf{Root cause.} A new RES series appeared in the data:
\texttt{wind\_off\_fcst\_mw} (Polish offshore wind). The first Baltic turbines
produced their first MWh on 2026-07-01. The training period (2023--2025)
had all zeros for this column. After the z-score normalisation step,
the training standard deviation was $0$, clamped to $10^{-6}$. A validation
value of $19$\,MW became $19 / 10^{-6} = 19{,}000{,}000$ in standardised
units. This created astronomically large network inputs.

\textbf{Why PatchTST failed but TFT did not.} TFT uses gated activations
(sigmoid gates in GRN and LSTM). Sigmoid saturates at extreme inputs:
$\sigma(10^7) \approx 1$. This acts as a hard clip. PatchTST uses linear
projection followed by GELU: $\text{GELU}(x) \approx x$ for large $x$. No
saturation, so the extreme input amplified through all layers.

\textbf{Fix.} In \texttt{standardize\_covariates}: compute the raw std before
clamping. Where raw std $< 10^{-4}$ (constant in training), zero that column
in all sets after standardisation. The zero-variance column carries no training
signal anyway; zeroing it prevents OOD values from exploding.

\begin{verbatim}
f_sd_raw = train.fut[:, :, :-n_tail].std(dim=(0, 1))
f_zero_mask = f_sd_raw < 1e-4
# after z-scoring:
s.fut[:, :, :-n_tail].masked_fill_(f_zero_mask, 0.0)
\end{verbatim}

\textbf{Result.} After fix: val pinball drops to $0.143$ for
\texttt{patch12\_s6\_ctx672} --- the expected range.

\textbf{Production lesson.} In live energy systems, new RES types enter
service every year (offshore wind, large-scale battery, hydrogen). A column
that is zero today will become non-zero next year. Any z-score pipeline
that relies on \texttt{clamp\_min(eps)} without a zero-variance guard is a
ticking deployment failure.

\subsection{Interview lines}

\begin{itemize}
  \item ``PatchTST reduces attention from $O(T^2)$ to $O((T/P)^2)$
        by treating price subsequences as tokens. This lets us use
        84-day context at the compute cost of a 4-day raw attention window.''
  \item ``Channel independence is a tradeoff: we lose solar $\times$ demand
        interactions, but we gain 20$\times$ fewer parameters --- which
        matters more when training samples are scarce.''
  \item ``The expected sweet spot is 12-hour patches: they preserve the
        morning/evening price regime split while still giving 5$\times$
        cheaper attention than raw hourly encoding.''
  \item ``We hit a real production bug: Baltic offshore wind came online in
        July 2026. The new column was all-zeros in training, so z-scoring
        divided by $10^{-6}$ and val inputs became $10^7$-scale. PatchTST
        exploded; TFT survived because sigmoid gates saturated. The fix is
        to detect zero-variance training columns and zero them rather than
        dividing by the clamp minimum.''
\end{itemize}
